% Part:sets-functions-relations
% Chapter: size-of-sets
% Section: non-enumerability-alt

\documentclass[../../../include/open-logic-section]{subfiles}

\begin{document}

\olfileid{sfr}{siz}{nen-alt}

\olsection{\printtoken{S}{nonenumerable} verzamelingen}

\begin{editorial}
  Deze paragraaf bewijst de overaftelbaarheid van $\Bin^\omega$ en
  $\Pow{\Nat}$ met de definities uit \olref[enm-alt]{sec}. Daarin wordt
  dus een bijectie met~$\Nat$ vereist in plaats van een surjectie
  vanuit $\PosInt$.
\end{editorial}

\begin{explain}
De verzameling $\Nat$ van natuurlijke getallen is oneindig. Zij is
bovendien triviaal !!{enumerable}. Opmerkelijk is echter dat er
\emph{!!{nonenumerable}} verzamelingen bestaan, dat wil zeggen
verzamelingen die niet !!{enumerable} zijn (zie
\olref[sfr][siz][enm-alt]{defn:enumerable}).

Dat kan verrassend zijn. Dat $A$ !!{nonenumerable} is, betekent immers
dat er \emph{geen} !!{bijection} $f \colon \Nat \to A$ bestaat. Er is
dus geen functie die de oneindig vele !!{element}s van~$\Nat$ naar~$A$
afbeeldt en daarbij heel~$A$ bestrijkt. Als $A$ !!{nonenumerable} is,
heeft $A$ in die zin ``meer'' !!{element}s dan er natuurlijke getallen
zijn.

Om te bewijzen dat een verzameling !!{nonenumerable} is, moet men
aantonen dat zo'n !!{bijection} niet kan bestaan. Dit kan het best door
te laten zien dat iedere poging de !!{element}s van~$A$ op te sommen
ten minste één !!{element} overslaat; daaruit volgt dat geen functie
$f\colon \Nat \to A$ !!{surjective} is. Een algemene strategie
hiervoor is Cantors \emph{diagonaalmethode}. Gegeven een lijst van
!!{element}s van $A$, bijvoorbeeld $x_1$, $x_2$, \dots, construeren we
een ander !!{element} van~$A$ dat door de wijze waarop het is
geconstrueerd onmogelijk op die lijst kan staan.

Dit wordt het duidelijkst aan de hand van een voorbeeld. Ons eerste
voorbeeld is de verzameling~$\Bin^\omega$ van alle oneindige rijen van
$0$'en en $1$'en. (Het symbool `$\Bin$' staat voor binair en kan
eenvoudig worden opgevat als de verzameling met twee elementen
$\{0,1\}$.)\oliflabeldef{sfr:card-arithmetic:card-opps:sec}{\footnote{Strikt
genomen zouden we $\Bin^\omega$ moeten vastleggen als de verzameling
van alle $\omega$-rijen van $0$'en en $1$'en, dat wil zeggen als de
verzameling $\funfromto{\omega}{\{0,1\}}$. Wat dit betekent, wordt
echter pas duidelijk in
\olref[sfr][card-arithmetic][card-opps]{sec}.}{} Voorlopig zal deze
enigszins onnauwkeurige formulering echter niet tot verwarring leiden.}
\end{explain}

\begin{thm}
\ollabel{thm:nonenum-bin-omega}
De verzameling $\Bin^\omega$~is !!{nonenumerable}.
\end{thm}

\begin{proof}
Beschouw een willekeurige opsomming van een deelverzameling van
$\Bin^\omega$. We hebben dan een lijst $s_{0}$, $s_{1}$, $s_{2}$,
\dots{}, waarbij iedere $s_n$ een oneindige rij van $0$'en en $1$'en is.
Laat $s_n(m)$ het cijfer zijn dat in de $n$-de rij op de $m$-de plaats
staat.
We kunnen de lijst nu opvatten als een tabel waarin $s_n(m)$ op de
$n$-de rij en in de $m$-de kolom staat:
\[
\begin{array}{c|c|c|c|c|c}
& 0 & 1 & 2 & 3 & \dots \\\hline
0 & \mathbf{s_{0}(0)} & s_{0}(1) & s_{0}(2) & s_0(3) & \dots \\\hline
1 & s_{1}(0)& \mathbf{s_{1}(1)} & s_1(2) & s_1(3) & \dots \\\hline
2 & s_{2}(0)& s_{2}(1) & \mathbf{s_2(2)} & s_2(3) & \dots \\\hline
3 & s_{3}(0)& s_{3}(1) & s_3(2) & \mathbf{s_3(3)} & \dots \\\hline
\vdots & \vdots & \vdots & \vdots & \vdots & \mathbf{\ddots}
\end{array}
\]
We construeren nu een oneindige rij~$d$ van $0$'en en $1$'en die niet
in deze lijst staat. Daartoe leggen we elk van haar cijfers vast: we
bepalen $d(n)$ voor alle $n \in \Nat$. Intuïtief lezen we de diagonaal
van de bovenstaande tabel van boven naar beneden---vandaar de naam
``diagonaalmethode''---en verwisselen we $1$ en $0$: waar de diagonaal
een $1$ bevat, kiezen we $0$, en omgekeerd. Formeler definiëren we
$d(n)$ als $0$ of $1$,
al naargelang het $n$-de !!{element} van de diagonaal, $s_n(n)$, gelijk
is aan $1$ of $0$:
\[
d(n) =
\begin{cases}
1 & \text{als $s_{n}(n) = 0$}\\
0 & \text{als $s_{n}(n) = 1$}
\end{cases}
\]
Er geldt duidelijk $d \in \Bin^\omega$, want het is een oneindige rij
van $0$'en en $1$'en. We hebben $d$ echter zo geconstrueerd dat
$d(n) \neq s_n(n)$ voor iedere $n \in \Nat$. De rij~$d$ verschilt dus
van~$s_n$ op haar $n$-de plaats. Bijgevolg geldt $d \neq s_n$ voor
iedere $n\in \Nat$, zodat $d$ niet kan voorkomen in de lijst $s_0$,
$s_1$, $s_2$, \dots

We hebben aangetoond dat een willekeurige opsomming van een
deelverzameling van $\Bin^\omega$ een !!{element} van $\Bin^\omega$
overslaat. Er bestaat dus geen opsomming van de verzameling
$\Bin^\omega$; met andere woorden, $\Bin^\omega$ is
!!{nonenumerable}.
\end{proof}

\begin{explain}
Deze bewijsmethode heet een ``diagonaalargument'', omdat de diagonaal van
de tabel wordt gebruikt om~$d$ te definiëren. Voor een diagonaalargument
is echter niet noodzakelijk een tabel aanwezig. Met een soortgelijk idee
kunnen we namelijk aantonen dat een verzameling !!{nonenumerable} is,
ook als er geen tabel en geen werkelijke diagonaal aan te pas komt. Het
volgende resultaat laat zien hoe.
\end{explain}

\begin{thm}
\ollabel{thm:nonenum-pownat}
$\Pow{\Nat}$ is niet !!{enumerable}.
\end{thm}

\begin{proof}
We gaan op dezelfde wijze te werk en tonen aan dat iedere lijst van
deelverzamelingen van~$\Nat$ een deelverzameling van $\Nat$ overslaat.
Stel dus dat we een lijst $N_0, N_1, N_2, \ldots$ van
deelverzamelingen van $\Nat$ hebben. We definiëren een verzameling~$D$
als volgt: $n \in D$ precies dan als $n \notin N_{n}$:
\[
D = \Setabs{n \in \Nat}{n \notin N_n}
\]
Er geldt duidelijk $D\subseteq \Nat$. De verzameling~$D$ kan echter
niet op de lijst staan. Volgens de constructie geldt immers $n \in D$
precies dan als $n\notin N_n$, zodat $D \neq N_n$ voor iedere
$n \in \Nat$.
\end{proof}

\begin{explain}
In het voorgaande bewijs werd geen diagonaal genoemd. Toch kan men er
een diagonaal in zien door de situatie als volgt weer te geven. Schrijf
de verzamelingen $N_0$, $N_1$, \dots, in een tabel en geef $N_n$ op de
$n$-de rij weer door $m$ in de $m$-de kolom te schrijven precies dan
als $m \in N_n$. Stel bijvoorbeeld dat de eerste vier verzamelingen op
de lijst $\{0,1,2,\dots\}$, $\{1, 3, 5, \dots\}$, $\{0,1,4\}$ en
$\{2,3,4,\dots\}$ zijn. De tabel begint dan als volgt:
\[
\begin{array}{r@{}rrrrrrr}
  N_0 = \{ & \mathbf{0}, & 1, & 2, & & & & \dots\}\\
  N_1 = \{ &  & \mathbf{1}, &  & 3, &  & 5, & \dots\}\\
  N_2 = \{ & 0, & 1, &  &  & 4\phantom{,} &  & \}\\
  N_3 = \{ &  &  & 2, & \mathbf{3}, & 4, & & \dots\}\\
  &\vdots & & & & & & \ddots\phantom{\}}
  \end{array}
\]
De verzameling~$D$ ontstaat nu door langs de diagonaal naar beneden te
gaan en $n \in D$ te nemen precies dan als $n$ \emph{niet} op de
diagonaal staat. In het bovenstaande geval laten we dus $0$ en $1$
weg, nemen we~$2$ op, laten we~$3$ weg, enzovoort.
\end{explain}

\begin{prob}
Toon met een expliciet diagonaalargument aan dat de verzameling van
alle functies $f \colon \Nat \to \Nat$ !!{nonenumerable} is. Stel dus
dat $f_1$, $f_2$, \dots, een lijst van functies is en dat iedere
$f_i\colon \Nat \to \Nat$ een functie is. Toon aan dat er dan een
functie $g \colon \Nat \to \Nat$ bestaat die niet op deze lijst staat.
\end{prob}
\end{document}